\section{Linearisation about the upright equilibrium}\label{sec:lin}

The equations of motion \cref{eq:eom-x,eq:eom-th} are nonlinear in $\theta$ and
$\dot\theta$. Near the upright equilibrium they can be replaced by a linear
time-invariant model
\begin{equation}
  \dot{\bm s} = \bm A\,\bm s + \bm B_{\mathrm{lin}}\,\tau ,
  \qquad \bm s = \begin{bmatrix} x & \theta & \dot x & \dot\theta \end{bmatrix}^{\mathsf T},
  \label{eq:lti}
\end{equation}
with constant matrices $\bm A\in\mathbb R^{4\times4}$ and
$\bm B_{\mathrm{lin}}\in\mathbb R^{4\times1}$. A model of this form is
introduced because its behaviour is fully characterised by the two matrices:
stability follows from the eigenvalues of $\bm A$, controllability from a rank
condition, and the standard methods of linear control design apply to it
directly. None of these tools applies to the nonlinear model.

\subsection{First-order approximation}

The equilibrium is $\theta = 0$, $\dot x = \dot\theta = 0$, $\tau = 0$, for any
$x$ (\cref{sec:model}). Near it $\theta$, $\dot x$, $\dot\theta$ and $\tau$ are
small quantities, and each term of the equations is kept only up to first order
in them. The trigonometric functions are replaced by their first-order Taylor
expansion about $\theta_0 = 0$,
\begin{equation}
  f(\theta) \approx f(0) + f'(0)\,\theta ,
\end{equation}
which gives
\begin{equation}
  \sin\theta \approx \sin 0 + \cos 0\cdot\theta = \theta ,
  \qquad
  \cos\theta \approx \cos 0 - \sin 0\cdot\theta = 1 .
\end{equation}
Term by term:
\begin{itemize}
  \item $\mb l\cos\theta\,\ddot\theta \approx \mb l\,\ddot\theta$: the
        correction to $\cos\theta \approx 1$ is of order $\theta^2$ and
        multiplies $\ddot\theta$, so it is of third order;
  \item $-\mb l\sin\theta\,\dot\theta^2 \approx -\mb l\,\theta\,\dot\theta^2$:
        product of three small quantities, third order, neglected;
  \item $-\mb g l\sin\theta \approx -\mb g l\,\theta$: first order, kept.
\end{itemize}
With $b_0 = \mb l$, the value of $b = \mb l\cos\theta$ at $\theta = 0$, the
linearised equations of motion are
\begin{equation}
  \underbrace{\begin{bmatrix} a & b_0 \\ b_0 & c \end{bmatrix}}_{\M_0}
  \begin{bmatrix} \ddot x \\ \ddot\theta \end{bmatrix}
  + \begin{bmatrix} 0 \\ -\mb g l\,\theta \end{bmatrix}
  = \begin{bmatrix} 1/r \\ -1 \end{bmatrix}\tau ,
  \label{eq:lin-eom}
\end{equation}
with $a$ and $c$ from \cref{eq:abc}.

\subsection{State-space matrices}

Moving the gravity term to the right-hand side of \cref{eq:lin-eom},
\begin{equation}
  \M_0 \begin{bmatrix} \ddot x \\ \ddot\theta \end{bmatrix}
  = \begin{bmatrix} \tau/r \\ -\tau + \mb g l\,\theta \end{bmatrix} .
\end{equation}
The inverse of $\M_0$ is
\begin{equation}
  \M_0^{-1} = \frac{1}{D}\begin{bmatrix} c & -b_0 \\ -b_0 & a \end{bmatrix},
  \qquad D = \det\M_0 = ac - b_0^2 > 0 ,
\end{equation}
so that
\begin{align}
  \ddot x &= \frac{1}{D}\left[c\,\frac{\tau}{r} - b_0\left(-\tau + \mb g l\,\theta\right)\right]
           = \frac{1}{D}\left[-b_0\,\mb g l\,\theta + \left(\frac{c}{r} + b_0\right)\tau\right],
  \label{eq:lin-xdd}\\
  \ddot\theta &= \frac{1}{D}\left[-b_0\,\frac{\tau}{r} + a\left(-\tau + \mb g l\,\theta\right)\right]
           = \frac{1}{D}\left[a\,\mb g l\,\theta - \left(a + \frac{b_0}{r}\right)\tau\right].
  \label{eq:lin-thdd}
\end{align}
The first two components of $\dot{\bm s}$ are $\dot x$ and $\dot\theta$, the
last two are \cref{eq:lin-xdd,eq:lin-thdd}. Writing $\alpha = -b_0\mb g l/D$
and $\beta = a\mb g l/D$,
\begin{equation}
  \bm A = \begin{bmatrix}
    0 & 0 & 1 & 0 \\
    0 & 0 & 0 & 1 \\
    0 & \alpha & 0 & 0 \\
    0 & \beta & 0 & 0
  \end{bmatrix},
  \qquad
  \bm B_{\mathrm{lin}} = \begin{bmatrix} 0 \\ 0 \\ b_3 \\ b_4 \end{bmatrix}
  = \begin{bmatrix} 0 \\ 0 \\ (c/r + b_0)/D \\ -(a + b_0/r)/D \end{bmatrix} .
  \label{eq:AB}
\end{equation}
With the parameters of \cref{tab:nominal}: $\alpha = \SI{-10.34}{s^{-2}}$,
$\beta = \SI{176.48}{s^{-2}}$, $b_3 = \SI{123.50}{kg^{-1}.m^{-1}}$,
$b_4 = \SI{-1309.37}{kg^{-1}.m^{-2}}$. The same matrices were obtained
numerically, by central finite differences of the nonlinear $\bm f(\bm s,\tau)$
of \cref{eq:forward} about the equilibrium with step \num{1e-6}; the largest
difference between the two is \num{2.2e-10}.

\subsection{Eigenvalues}

The eigenvalues of $\bm A$ are the roots of $\det(\lambda\bm I - \bm A) = 0$:
\begin{equation}
  \det\begin{bmatrix}
    \lambda & 0 & -1 & 0 \\
    0 & \lambda & 0 & -1 \\
    0 & -\alpha & \lambda & 0 \\
    0 & -\beta & 0 & \lambda
  \end{bmatrix} .
\end{equation}
The first column has a single non-zero entry, so expanding along it,
\begin{equation}
  \det(\lambda\bm I - \bm A)
  = \lambda \det\begin{bmatrix}
    \lambda & 0 & -1 \\
    -\alpha & \lambda & 0 \\
    -\beta & 0 & \lambda
  \end{bmatrix}
  = \lambda\left[\lambda\cdot\lambda^2 - 1\cdot\left(0 + \lambda\beta\right)\right]
  = \lambda^2\left(\lambda^2 - \beta\right) .
\end{equation}
The eigenvalues are therefore
\begin{equation}
  \lambda_{1,2} = 0, \qquad \lambda_{3,4} = \pm\sqrt{\beta} = \pm\sqrt{\frac{a\,\mb g l}{D}} = \pm\SI{13.28}{s^{-1}} .
\end{equation}
The double zero belongs to the position: with $\tau = 0$ nothing acts on $x$
other than the coupling with $\theta$, and the robot has no preferred position.
The positive eigenvalue makes the equilibrium unstable: a deviation in
$\theta$ grows as $e^{\sqrt\beta\,t}$, by a factor $e$ every
$1/\sqrt\beta = \SI{75}{ms}$. The coefficient $\alpha$ does not affect the
eigenvalues: the motion of the wheels does not act back on $\theta$.

\paragraph{Comparison with the nonlinear model.} With $\tau = 0$,
\cref{eq:lin-thdd} reduces to $\ddot\theta = \beta\,\theta$, whose general
solution is $\theta(t) = C_1 e^{\sqrt\beta t} + C_2 e^{-\sqrt\beta t}$. From rest,
$\theta(0) = \theta_0$ and $\dot\theta(0) = 0$ give $C_1 = C_2 = \theta_0/2$,
hence
\begin{equation}
  \theta(t) = \theta_0\cosh\!\left(\sqrt\beta\, t\right) .
\end{equation}
\Cref{tab:lin-vs-nl} compares this solution with the nonlinear simulation of
\cref{sec:sim} (RK4, $\Delta t = \SI{1}{ms}$), both from $\theta_0 = 3^\circ$.

\begin{table}[h]
\centering
\caption{Free fall from $\theta_0 = 3^\circ$: linearised and nonlinear model.}
\label{tab:lin-vs-nl}
\begin{tabular}{@{}ccc@{}}
\toprule
$t$ [s] & $\theta$, linear [deg] & $\theta$, nonlinear [deg] \\
\midrule
0.05 & 3.7  & 3.7  \\
0.10 & 6.1  & 6.0  \\
0.15 & 11.2 & 11.1 \\
0.20 & 21.5 & 20.7 \\
0.25 & 41.6 & 37.2 \\
0.30 & 80.7 & 62.0 \\
\bottomrule
\end{tabular}
\end{table}

\paragraph{Limits of the linear model.} The two models agree within
$1^\circ$ up to about $20^\circ$ and diverge beyond. The linear model uses
$\theta$ in place of $\sin\theta$; since $\theta > \sin\theta$ for
$\theta > 0$, it overestimates the gravitational torque and predicts a faster
fall. The relative error of the approximation is $1\%$ at $14^\circ$ and
$4.7\%$ at $30^\circ$, and the neglected terms (such as
$\theta\dot\theta^2$) also grow with the speed of the motion. The linear model
is therefore a valid description only in a neighbourhood of the equilibrium,
and results obtained from it must be checked on the nonlinear model.

\subsection{Controllability}

The system \cref{eq:lti} is controllable, i.e.\ the single input $\tau$ can
bring the state from any initial value to any final value in finite time, if
and only if the controllability matrix
\begin{equation}
  \bm C = \begin{bmatrix} \bm B_{\mathrm{lin}} & \bm A\bm B_{\mathrm{lin}} & \bm A^2\bm B_{\mathrm{lin}} & \bm A^3\bm B_{\mathrm{lin}} \end{bmatrix}
\end{equation}
has rank 4. From \cref{eq:AB},
\begin{equation}
  \bm B_{\mathrm{lin}} = \begin{bmatrix} 0 \\ 0 \\ b_3 \\ b_4 \end{bmatrix},\quad
  \bm A\bm B_{\mathrm{lin}} = \begin{bmatrix} b_3 \\ b_4 \\ 0 \\ 0 \end{bmatrix},\quad
  \bm A^2\bm B_{\mathrm{lin}} = \begin{bmatrix} 0 \\ 0 \\ \alpha b_4 \\ \beta b_4 \end{bmatrix},\quad
  \bm A^3\bm B_{\mathrm{lin}} = \begin{bmatrix} \alpha b_4 \\ \beta b_4 \\ 0 \\ 0 \end{bmatrix}.
\end{equation}
Reordering the columns as (2, 4, 1, 3), which changes at most the sign of the
determinant, gives a block-diagonal matrix with two equal blocks
\begin{equation}
  \bm G = \begin{bmatrix} b_3 & \alpha b_4 \\ b_4 & \beta b_4 \end{bmatrix},
  \qquad |\det\bm C| = (\det\bm G)^2 ,
\end{equation}
with
\begin{equation}
\begin{aligned}
  \det\bm G &= b_4\left(\beta b_3 - \alpha b_4\right)
  = b_4\,\frac{\mb g l}{D^2}\left[a\left(\frac{c}{r} + b_0\right) - b_0\left(a + \frac{b_0}{r}\right)\right] \\
  &= b_4\,\frac{\mb g l}{D^2}\,\frac{ac - b_0^2}{r}
  = b_4\,\frac{\mb g l}{r D} .
\end{aligned}
\end{equation}
Since $b_4 \neq 0$ and $D > 0$, $\det\bm C \neq 0$ for every value of the
parameters: the linearised system is controllable, although it has two degrees
of freedom and a single input. The rank was also confirmed numerically
(\texttt{numpy.linalg.matrix\_rank}) with the parameters of \cref{tab:nominal}.
